\chapter{Application to the fundamental group} \label{X}

\renewcommand{\theequation}{\arabic{equation}}

Throughout this expos\'e, one will denote by $X$ a locally noetherian
prescheme, by~$Y$ a closed subset of $X$, by $U$ a variable open
neighbourhood of $Y$ in $X$ and by~$\hat{X}$ the formal completion of
$X$ along $Y$ (\EGA I 10.8). For every prescheme $Z$, one will denote
by $\Et(Z)$ the category of \'etale coverings of $Z$, and by
$\LL(Z)$\label{pX.1} the category of locally free coherent Modules on
$Z$.

%
\section{Comparison of $\Et(\hat{X})$ and of $\Et(Y)$} \label{X.1}

Let $\Ii$ be an ideal of definition of $Y$ in $X$. Set, for every
$n\in\NN$, $Y_n=(Y, (\OX/\Ii^{n+1})_{|Y})$. The $Y_n$ form an inductive
system of ordinary preschemes, or also of formal preschemes, by
endowing the structure sheaves with the discrete topology. One knows
(\EGA I 10.6.2) that $\hat{X}$ is the inductive limit, in the category
of formal preschemes, of the inductive system of the $Y_n$. One also
knows (\EGA I 10.13) that to give a formal $\hat{X}$-prescheme of
\emph{finite type} $R$ is to give an inductive system of ordinary
$Y_n$-preschemes $R_n$ of \emph{finite type}, such that
$R_n\simeq(R_{n+1})\times_{(Y_{n+1})}(Y_n)$. Moreover, for $R$ to be an
\'etale covering of $\hat{X}$, it is necessary and sufficient that for
every~$n$, $R_n$ be an \'etale covering of $Y_n$. This being said, it
is easy to see that nilpotent elements do not matter for \'etale
coverings (\SGA 1~8.3), i.e.\ that the change of base functor:
\[
\Et(Y_{n+1}) \to \Et(Y_n)
\]
is an equivalence of categories for every $n\in\NN$. Hence:

\begin{proposition} \label{X.1.1}
With the notations introduced above, the natural functor
$\Et(\hat{X})\to\Et(Y)$ is an equivalence of categories
(\cf \textup{\SGA 1~8.4}).
\end{proposition}

\section{Comparison of $\Et(Y)$ and $\Et(U)$, for $U$ variable} \label{X.2}

We shall introduce two conditions from which the announced comparison
theorem will follow easily. Let $X$ be a locally noetherian prescheme
and let $Y$ be a closed subset of $X$. One says that the pair $(X, Y)$
satisfies the \emph{Lefschetz condition}, which one writes
$\Lef(X, Y)$,\label{pX.2} if, for every open $U$ of $X$ containing $Y$
and every locally free coherent sheaf $E$ on $U$, the natural
homomorphism
\[
\Gamma(U, E) \to \Gamma(\hat{X}, \hat{E})
\]
is an isomorphism.

One says that the pair $(X, Y)$ satisfies the \emph{effective Lefschetz
condition}, which one writes $\Leff(X, Y)$, if one has $\Lef(X, Y)$ and
if moreover, for every locally free coherent sheaf~$\Ee$ on $\hat{X}$,
there exist an open neighbourhood $U$ of $Y$ and a locally free
coherent sheaf~$E$ on $U$ and an isomorphism $\hat{E}\simeq\Ee$.

These conditions are satisfied in two important examples:

\refstepcounter{subsection}\label{X.2.1}
\begin{enonce*}[remark]{Example 2.1\ndemark}
\ndetext{\label{noteX.2.1}one can slightly improve (i): see Mme Raynaud
(Raynaud~M., {\og Th\'eor\`emes de Lefschetz en cohomologie des
faisceaux coh\'erents et en cohomologie \'etale. Application au groupe
fondamental\fg}, \emph{Ann. Sci. \'Ec. Norm. Sup. (4)} \textbf{7}
(1974), p\ptbl 29--52, corollaries I.1.4 and I.5); condition (ii) can
be improved so as to get rid of the depth conditions along $Y$ (see
theorem 3.3 of \loccit for a precise statement). The proof of this last
point is very technical, the article above moreover giving only
indications of the proof, referring to a detailed earlier version,
which appeared in the Bulletin de la Soci\'et\'e math\'ematique de
France.} Let $A$ be a noetherian ring and let $t\in \rr(A)$ be an
$A$-regular element belonging to the radical $\rr(A)$ of $A$. Suppose
that $A$ is a quotient of a regular local ring and that $A$ is complete
for the $t$-adic topology (for example $A$ complete for the
$\rr(A)$-adic topology). Set $X'=\Spec(A)$ and $Y'=V(t)$, and moreover
set $x=\rr(A)$ and $X=X'-\{x\}$, $Y=Y'-\{x\}$. Thus $X$ is open in $X'$
and $Y=X\cap Y'$. Then:
\begin{enumeratei}
\item
If, for every prime ideal $\pp$ of $A$ such that $\dim A/\pp=1$
(\ie for every closed point of $X$) one has $\prof A_{\pp}\geq 2$,
then one has $\Lef(X, Y)$;

\item
if moreover, for every prime ideal $\pp$ of $A$ such that $t\in\pp$
and $\dim A/\pp=1$ (\ie for every closed point of $Y$), one has
$\prof A_{\pp}\geq 3$, then one has $\Leff(X, Y)$.
\end{enumeratei}
\end{enonce*}

Let us first show that for every open neighbourhood $U$ of $Y$ in $X$,
the complement of $U$ in $X$ is a union of finitely many closed points
(in $X$). Let us remark that $U$ is open in $X$, hence in $X'$, hence
$Z'=X'-U$ is closed. Let $I$ be an ideal of definition of $Z'$; it
suffices to prove that $A/I$ is of dimension $1$. Now $Z'\cap Y'=\{x\}$
hence $A/(I+(t))$ is artinian, whence the conclusion thanks to the
``Hauptidealsatz''.

The \emph{first hypothesis is equivalent to}: ``for every prime ideal
$\pp$ of $A$, $\pp\neq \rr(A)$, one has
$\prof A_\pp\geq3-\dim A/\pp$''. Indeed $A$ is a quotient of a regular
ring, one can therefore apply \Ref{VIII}~\Ref{VIII.2.3} to the
prescheme $X'$, to the closed subset $\{x\}$ and to the coherent sheaf
$\Oo_{X'}$, observing that, $c(\pp)=\dim(A/\pp)$ for
$\pp\in U=X'-\{x\}$ (because $x$ is the closed point of $X'$).

Let $U$ be an open neighbourhood of $Y$ in $X$ and $E$ a locally free
$\Oo_U$-module. Set $Z=X-U$ and let $u\colon U\to X$ be the canonical
immersion. One will first prove that $u_\ast(E)$ is a \emph{coherent}
$\OX$-Module, or, what amounts to the same, that $\SheafH^i_Z(E')$ is
coherent for $i=0, 1$, where $E'$ is a coherent extension of $E$ to
$X$. For this one applies theorem \Ref{VIII}~\Ref{VIII.2.1} to the
prescheme $X$, to the closed subset $Z$ and to the coherent sheaf
$E'$. It suffices to verify that for every point $p\in U$ such that
$c(p)=1$, one has $\prof E'_p\geq 1$, where one has set
\[
c(p)=\codim (\overline{\{p\}}\cap Z, \overline{\{p\}}).
\]
Now if $p\in U$ and if $c(p)=1$, denoting by $\pp$ the ideal of $A$
corresponding to $p$, one sees that $\dim A/\pp=2$, because the
complement of $U$ is a union of finitely many closed points and $A$ is
a quotient of a regular ring. Moreover, $E$ is locally free, hence, for
every $p\in\Supp E$, one has $\prof E_p=\prof\Oo_{U, p}$. Finally, if
$p\in U$ and if $c(p)=1$, one has
\[
\prof E'_p = \prof E_p = \prof\Oo_{U, p} = \prof A_\pp \geq 3-2=1.
\]

We must now prove that the natural homomorphism
\begin{equation} \label{eq:X.1} {\Gamma(U, E) \to \Gamma(\hat{X}, \hat{E})}
\end{equation}
is an isomorphism. Setting then $\overline{\!E}=u_\ast(E)$, one notes
that $\overline{\!E}$ is coherent and of depth $\geq 2$ at all the
closed points of $X$. It follows that $\R^i f_\ast(\overline{\!E})$ is
coherent for $i=0, 1$, where $f\colon X\to X'$ denotes the canonical
immersion of $X$ in $X'=\Spec(A)$ (\Exp \Ref{VIII}). One then applies
(\Ref{IX}~\Ref{IX.1.5}), and one concludes that
\begin{equation} \label{eq:X.2}
\Gamma(U, \overline{\!E}) \to \Gamma(\hat{X}, \widehat{\overline{\!E}})
\end{equation}
is an isomorphism, because $A$ is complete for the $t$-adic topology.

One has a commutative diagram:
\[
\UseTips \newdir{ >}{!/-5pt/\dir{>}} \xymatrix@=5mm{ \Gamma(X, \overline{\!E})\ar[rr]^{\simeq}\ar[dr]_{\simeq} && \Gamma(U, E)\ar[dl]\\
&\Gamma(\hat{X}, \hat{E})}
\]
whence the conclusion.

Let now $\Ee$ be a locally free coherent sheaf on $\hat{X}$. If one has
proved that~$\Ee$ is algebraizable, \ie is isomorphic to the formal
completion of a coherent $\OX$-Module~$E$, it is easy to see that $E$
is locally free in a neighbourhood of $Y$, hence to prove
$\Leff (X, Y)$. Let $\widehat{X'}$ be the formal spectrum of $A$ for
the $t$-adic topology, which is identified with the formal completion
of $X'$ along $Y'$. Let us denote by $f$ the canonical immersion of~$X$
in $X'$, by $f'$ the canonical immersion of $Y$ in $Y'$, and by
$\hat{f}$ the morphism deduced by passage to the completions. For
$\Ee$ to be algebraizable it suffices that $\hat{f}_\ast(\Ee)$ be a
coherent $\Oo_{\hat{X}}$-Module, because $A$ is complete for the
$t$-adic topology. Let $\Ii = t\Oo_{\hat{X}}$, it is an ideal of
definition of $\hat{X}$.

For every $n\geq 0$, set $\E_n=\Ee/\Ii^{n+1}\Ee$. At every closed point
$y\in Y$, the depth of $\E_0$ is $\geq 2$; indeed $t$ is an $A$-regular
element hence $\prof\Oo_{Y_0, y}=\prof\Oo_{X, y}-1\geq 2$. One concludes
that $\hat{f}_\ast(\Ee)$ is coherent (\Ref{IX}~\Ref{IX.2.3}). qed

\refstepcounter{subsection}\label{X.2.2}
\begin{enonce*}[remark]{Example 2.2}
{\normalfont (Will allow one to compare the fundamental group of a projective variety and of a hyperplane section)}
Let $K$ be a field and let $X$ be a proper $K$-prescheme. Let $\Ll$ be
an ample invertible $\OX$-Module. Let $t \in \Gamma(X, \Ll)$ be an
$\OX$-regular element, which means that, for every open $U$ and every
isomorphism $u\colon \Ll_{|U}\to\Oo_U$, \/$u(t)$ is not a zero-divisor
in $\Oo_U$ (a condition which does not depend on $u$). Let $Y = V(t)$
be the subscheme of $X$ with equation $t = 0$.\nde{condition (ii) is
superfluous, see the note on page~\pageref{X.2.1}.} Then:
\begin{enumeratei}
\item
If, for every closed point $x$ in $X$, one has
$\prof\Oo_{X, x}\geq 2$, one has $\Lef (X, Y)$;

\item
if moreover, for every closed point $y\in Y$, one has
$\prof\Oo_{X, y}\geq 3$, one has $\Leff (X, Y)$.
\end{enumeratei}
\end{enonce*}

This example will be treated in detail in \Exp \Ref{XII}.

Let $S$ be a prescheme, one knows (\EGA II 6.1.2) that the functor
which, to every finite flat covering $r\colon R\to S$, associates the
$\OX$-Algebra $r_\ast(\Oo_R)$, induces an equivalence between the
category of finite flat coverings of $S$ and the category of coherent
and locally free $\OX$-Algebras. Let $U$ be an open neighbourhood of
$Y$, and let $r\colon R\to U$ be a finite flat covering of $U$. Let
$\hat{R}$ be the finite flat covering of $\hat{X}$ deduced from it by
change of base. One has $\hat{r}_\ast(\Oo_{\hat{R}})\simeq\widehat{r_\ast(\Oo_R)}$.

\emph{Suppose then} $\Lef(X, Y)$. This implies that, for every $U$, the
inverse image functor:
\[
\LL(U) \to \LL(\hat{X})
\]
is \emph{fully faithful}. Indeed, let $E$ and $F$ be two locally free
coherent $\Oo_U$-Modules; $\SheafHom(E, F)$ is also coherent and
locally free; $\SheafHom(E, F)$ is also coherent and locally free. By
hypothesis, the natural map
\[
\Gamma_U (\SheafHom(E, F)) \to \Gamma_{\hat{X}}(\widehat{\SheafHom(E, F)})
\]
is an isomorphism, whence the conclusion, because $\SheafHom$ commutes
with $\, \, \hat{}\, \,$ since everything is locally free. Now
$\, \, \hat{}\, \,$ commutes with the tensor product, one deduces that
the functor which to every coherent and locally free $\Oo_U$-Algebra
$\Aa$ associates the $\Oo_{\hat{X}}$-Algebra $\hat{\Aa}$, is fully
faithful. Better still, if $E$ is a locally free coherent
$\Oo_U$-Module, there is a one-to-one correspondence between the
structures of commutative $\Oo_{\hat{X}}$-Algebra on $\hat{E}$.

\begin{proposition} \label{X.2.3}
Let $X$ be a locally noetherian prescheme and let $Y$ be a closed
subset of $X$. Let $\hat{X}$ be the formal completion of $X$ along
$Y$. For every open $U$ of $X$, $U\supset Y$, denote by $L_U$
(\resp $P_U$, \resp $E_U$) the functor which to every locally free
coherent $\Oo_U$-Module (\resp to every finite flat covering of $U$,
\resp to every \'etale covering of $U$) associates its inverse image
under $\hat{X} \to X$.
\begin{enumeratei}
\item
If one has $\Lef(X, Y)$, then for every open neighbourhood $U$ of $Y$,
the functors $L_U$, $P_U$ and $E_U$ are fully faithful.

\item
If one has $\Leff(X, Y)$, then for every locally free coherent
$\Oo_{\hat{X}}$-Module~$\Ee$ (\resp ...), there exist an open $U$ and a
locally free coherent $\Oo_U$-Module~$E$ (\resp ...), such that
$L_U(E) \simeq \Ee$ (\resp ...).
\end{enumeratei}
\end{proposition}

(i) Has been seen.

(ii) Follows from (i) and the hypothesis, at least for $L_U$ and
$P_U$. Moreover if $R$ is an \emph{\'etale} covering of $\hat{X}$,
there exist an open neighbourhood $U$ of $Y$ in $X$ and a finite flat
covering $R'$ of $U$ such that $\widehat{R'} \simeq R$. One deduces
from this a covering $R''$ of $Y$ which is \'etale by \Ref{X.1.1},
hence $R'$ is \'etale in a neighbourhood $U'$ of $Y$. qed

\begin{corollary} \label{X.2.4}
If one has $\Lef (X, Y)$, for a finite flat covering $R$ of an open
neighbourhood $U$ of $Y$ to be connected, it is necessary and
sufficient that $R \times_U \hat{X}$ be so. In particular, for $Y$ to
be connected, it is necessary and sufficient that the open
neighbourhood $U$ of $Y$ be so, or again that $X$ be so.
\end{corollary}

Indeed, for a locally ringed space $(X, \OX)$ to be connected, it is
necessary and sufficient that $\Gamma(X, \OX)$ not be a direct product
of two nonzero rings. Now one has
\[
\Gamma(U, r_\ast(\Oo_R)) \simeq \Gamma(\hat{X}, \hat{r}_\ast(\Oo_{\hat{R}}))
\]
by $\Lef(X, Y)$.

\begin{corollary} \label{X.2.5} If one has $\Lef (X, Y)$, then for every $U$, the functor
\[
\Et(U) \to \Et(Y)
\]
is fully faithful. If one has $\Leff(X, Y)$, then for every \'etale
covering $R$ of $Y$, there exist an open neighbourhood $U$ of $Y$ and a
covering $R'$ of $U$ such that $R' \times_U Y \simeq R$.
\end{corollary}

\refstepcounter{subsection}\label{X.2.6}
\begin{enonce*}{Corollary 2.6\ndemark}
If one has $\Lef(X, Y)$ and if $Y$ is connected, every open
neighbourhood $U$ of $Y$ is connected and the natural homomorphism
$\pi_1(Y)\to\pi_1(U)$ is surjective. If moreover one has
$\Leff(X, Y)$, the natural homomorphism
\[
\pi_1(Y) \to \varprojlim_{U}\, \pi_1(U)
\]
is an isomorphism. (N.B. One assumes chosen a ``base-point'' in $Y$,
which one also takes as base-point in $X$, for the definition of the
fundamental groups.)
\end{enonce*}

All this follows trivially from prop\ptbl \Ref{X.1.1} and from
prop\ptbl \Ref{X.2.3}.
\ndetext{together with~\Ref{X.3.3} and the criteria~\Ref{XII.2.4}
and~\Ref{XII.3.4}, one obtains the following relative Lefschetz
theorem. Let $f:X\to S$ be a projective and flat morphism of connected
noetherian schemes and let $D$ be an effective relative Cartier divisor
in $X$ which is relatively ample. If, for every $s\in S$, the depth of
$X_s$ at each closed point is $\geq 2$, then $D$ is connected and, for
every open $U$ of $X$ containing $D$, the arrow
$i_U:\pi_1(D)\to\pi_1(U)$ is surjective. If moreover, the depth of
$X_s$ along each closed point of $D_s$ is $\geq 3$ and if the local
rings of $X$ at its closed points are pure, for example complete
intersections (\cf \Ref{X}~\Ref{X.3.4}), then $i_X$ is an isomorphism.
\Cf Bost~J.-B., {\og Lefschetz theorem for Arithmetic Surfaces\fg},
\emph{Ann. Sci. \'Ec. Norm. Sup. (4)} \textbf{32} (1999),
p\ptbl 241--312, theorems 1.1 and 2.1. In the case where $X$ is simply
a smooth projective geometrically connected surface over a field, one
always has connectedness of $D$ and surjectivity of
$\pi_1(D)\to\pi_1(U)$ (where $U$ is an open containing~$D$) for $D$
merely nef of square $>0$ (\cf \loccit, theorem 2.3 and also theorem
2.4 for surfaces that are merely normal and complete). In the case of
an arithmetic surface (normal and quasi-projective) $X$ over a ring of
integers $\Oo_K$, Bost, improving results of Ihara (Ihara~Y.,
{\og Horizontal divisors on arithmetic surfaces associated with
Bely\u\i\ uniformizations\fg}, in \emph{The Grothendieck theory of
dessins d'enfants (Luminy, 1993)}, London Math. Soc. Lect. Note Series,
vol.~200, Cambridge Univ. Press, Cambridge, 1994, 245--254 or \loccit,
corollary 7.2), has shown that if a point $P\in X(\Oo_K)$, which plays
the role of the divisor $D$ in the geometric situation, satisfies
certain positivity conditions, then the arrow
$\pi_1(X)\to\pi_1(\Spec \Oo_K)$ deduced from the projection was
invertible, with inverse the arrow
$\pi_1(\Spec \Oo_K)\to\pi_1(X)$ deduced from $P$ (\loccit,
theorem 1.2).}

\section{Comparison of $\pi_1(X)$ and of $\pi_1(U)$} \label{X.3}

\begin{definition} \label{X.3.1}
Let $X$ be a prescheme and $Z$ a closed subset of $X$. Set $U=X-Z$.
One says that the pair $(X, Z)$ is pure if, for every open $V$ of $X$,
the functor
\[
\begin{aligned} \Et(V) &\to \Et(V\cap U) \\
V' &\mto V'\times_V(V\cap U)
\end{aligned}
\]
is an equivalence of categories\sfootnote{For a notion more
satisfactory in certain respects, \cf the comment \Ref{XIV}~\Ref{XIV.1.6}
d).}.
\end{definition}

\begin{definition} \label{X.3.2}
Let $A$ be a noetherian local ring. Set $X=\Spec A$. Let $\rr(A)$ be
the radical of $A$ and let $x=\rr(A)$ be the closed point of $X$. One
says that $A$ is pure if the pair $(X, \{x\})$ is so.
\end{definition}

We leave to the reader the task of not proving the following
proposition:

\begin{proposition} \label{X.3.3}
Let $X$ be a locally noetherian prescheme and let $Z$ be a closed
subset of $X$. For the pair $(X, Z)$ to be pure it is necessary and
sufficient that, for every $z\in Z$, the ring $\Oo_{X, z}$ be
pure\sfootnote{Compare the non-commutative case of
\Ref{XIV}~\Ref{XIV.1.8}, whose proof is essentially the same as that of
\Ref{X.3.3}.}.
\end{proposition}

This being said the following theorem is the essential result of this
number:

\refstepcounter{subsection}\label{X.3.4}
\begin{enonce*}{Theorem 3.4 (Purity theorem\ndemark)}
\ndetext{for the history of the methods employed, see the letter of
1~October 1961 from Grothendieck to Serre, \emph{Correspondance
Grothendieck-Serre}, edited by Pierre Colmez and Jean-Pierre Serre,
Documents Math\'ematiques, vol.~2, Soci\'et\'e Math\'ematique de
France, Paris, 2001.}
\textup{(i)}\enspace
A regular noetherian local ring of dimension $\geq 2$ is pure (purity
theorem of Zariski--Nagata).

\textup{(ii)}\enspace
A noetherian local ring of dimension $\geq 3$ which is a complete
intersection is pure.
\end{enonce*}

Let us recall that one says that a local ring is a \emph{complete
intersection} if there exists a \emph{regular} noetherian local ring
$B$ and a $B$-regular sequence $(t_1, \ldots, t_k)$ of elements of the
radical $\rr(B)$ of $B$ such that
\[
A \simeq B/(t_1, \ldots, t_k).
\]

In this connection let us remark that it would be less ambiguous to say
that $A$ is an \emph{absolute} complete intersection, as opposed to the
situation, which we have already encountered, where $X$ is a locally
noetherian prescheme (which need not be regular) and where $Y$ is a
closed subset of $X$, of which one says that it is ``set-theoretically
locally a complete intersection \emph{in $X$}''.

Let us first prove some lemmas.

\begin{lemma} \label{X.3.5}
Let $X$ be a locally noetherian prescheme and let $U$ be an open subset
of $X$. Set $Z=X-U$. Let $i\colon U\to X$ be the canonical immersion of
$U$ in $X$. The following conditions are equivalent:
\begin{enumeratei}
\item
For every open $V$ of $X$, if one sets $V'=V\cap U$, the functor
$F\mto F_{|V'}$ from the category of locally free coherent
$\Oo_V$-Modules to the category of locally free coherent
$\Oo_{{V'}}$-Modules is fully faithful;

\item
the natural homomorphism $\OX \to i_\ast(\Oo_{U})$ is an isomorphism;

\item
for every $z\in Z$, one has $\prof\Oo_{X, z}\geq 2$.
\end{enumeratei}
\end{lemma}

One has already seen (\Ref{III}~\Ref{III.3.3}) the equivalence of (ii)
and (iii). Let us show that (ii) implies (i). Let $F$ and $G$ be two
locally free coherent $\Oo_V$-Modules, $\SheafHom(F, G)$ is so as well,
hence $\SheafHom(F, G)\to i_\ast(\SheafHom(F_{|V'}, G_{|V'}))$ is an
isomorphism hence $\Hom(F, G)\simeq\Hom(F_{|V'}, G_{|V'})$. Conversely
one takes $F=G=\OX$ and one applies (i) to every open~$V$ of~$X$.

Here is a useful ``descent lemma'':

\begin{lemma} \label{X.3.6} Let $X$ be a locally noetherian prescheme and let $Z$ be a closed subset of $X$. Set $U=X-Z$. Suppose that the homomorphism $\OX \to i_\ast(\Oo_U)$ is an isomorphism. Let $f\colon X_1\to X$ be a faithfully flat and quasi-compact morphism. Set $Z_1=f^{-1}(Z)$. If the pair $(X_1, Z_1)$ is pure, so is $(X, Z)$.
\end{lemma}

Let us remark that the hypothesis $\OX\simeq i_\ast(\Oo_U)$ is
preserved by flat extension of the base, because $i$ is a quasi-compact
morphism and, in this case, direct image commutes with inverse image.
Now this hypothesis implies that the functor
\[
\Et(V) \to \Et(U\cap V)
\]
defined by
\[
V'\mto V'\times_V(V\cap U)
\]
is fully faithful, as is shown by the interpretation of an \'etale
covering in terms of a locally free coherent Algebra. It remains to
prove effectivity. One may for example introduce the square $X_2$ and
the cube $X_3$ of $X_1$ over $X$ and one remarks that a faithfully flat
and quasi-compact morphism is a morphism of \emph{universal effective
descent} for the fibered category of \'etale coverings, over the
category of preschemes. The conclusion is formal from
there\sfootnote{\Cf J.~Giraud, \emph{M\'ethode de la descente},
M\'emoire \numero 2 of the Bulletin de la Soci\'et\'e Math\'ematique de
France (1964).}.

\begin{remark} \label{X.3.7}
One has proved along the way that if $\OX\to i_\ast(\Oo_U)$ is an
isomorphism, $X$ is connected if and only if $U$ is so, and then
$\pi_1(U)\to\pi_1(X)$ is surjective.
\end{remark}

\begin{corollary} \label{X.3.8}
Let $A$ be a noetherian local ring. Suppose that $\prof A\geq 2$. Then
if $\hat{A}$ is pure, $A$ is pure.
\end{corollary}

Follows from lemma \Ref{X.3.5} and lemma \Ref{X.3.6}.

The following lemma is the essential point of the proof of the purity
theorem:

\begin{lemma} \label{X.3.9}
Let $A$ be a noetherian local ring and let $t\in\rr(A)$ be an
$A$-regular element. Suppose that $A$ is complete for the $t$-adic
topology and moreover a quotient of a regular local ring (for example
$A$ complete). Set $B=A/tA$.
\begin{enumeratei}
\item
If for every prime ideal $\pp$ of $A$ such that $\dim A/\pp=1$, one has
$\prof A_{\pp}\geq 2$, then $B$ pure implies $A$ pure.

\item
If for every prime ideal $\pp$ of $A$ such that $\dim A/\pp=1$, one has
$\prof A_{\pp}\geq 2$, if $A_\pp$ is pure when $t\notin\pp$, and
if\nde{this last condition can be improved, \cf the editor's
note~\eqref{noteX.2.1} on page~\pageref{noteX.2.1}.} $\prof A_\pp\geq3$
when $t\in\pp$, then $A$ pure implies $B$ pure.
\end{enumeratei}
\end{lemma}

Let $X'=\Spec(A)$ and let $Y'=V(t)$, which one identifies with the
spectrum of $B$. Let $x=\rr(A)$, and set $X=X'-\{x\}$ and
$Y=Y'-\{x\}=X\cap Y'$. Let us denote by $\widehat{X'}$ the formal
spectrum of $A$ for the $t$-adic topology, which is identified with the
formal completion of $X'$ along $Y'$.

Since $A$ is complete for the $t$-adic topology, one remarks that
$\Et(X')\to\Et(\widehat{X'})$ is an equivalence of categories. Likewise
$\Et(\widehat{X'})\to\Et(Y')$ by prop\ptbl \Ref{X.1.1}, hence
$\Et(X')\to\Et(Y')$ is an equivalence of categories.

Let us show (i). Consider the diagrams
\[
\UseTips \newdir{ >}{!/-5pt/\dir{>}} \xymatrix{ X' & X\ar[l] && \Et(X')\ar[r]^a \ar[d]_c & \Et(X)\ar[d]^b\\
Y'\ar[u] & Y\ar[l] \ar[u] && \Et(Y')\ar[r]^d & \Et(Y)}
\]
One has just seen that $c$ is an equivalence, $d$ is so as well by the
hypothesis that $B$ is pure and finally $b$ is fully faithful as one
has seen in example \Ref{X.2.1}, \cf \Ref{X.2.3} (i).

Let us show (ii). This time one assumes that $A$ is pure hence $a$ is
an equivalence; likewise $c$. Let us see that $b$ is an equivalence. By
example \Ref{X.2.1} one knows that one has $\Leff(X, Y)$, hence $b$ is
already fully faithful, let us prove that it is essentially
surjective. One uses \Ref{X.2.3} (ii) noting that, if $U$ is an open
neighbourhood of $Y$ in $X$, the complement of $U$ in $X$ is a union of
finitely many closed points; the pair $(X,X-U)$ is therefore pure by
prop\ptbl \Ref{X.3.3}, because at such a point $\pp$,
$\Oo_{X,\pp}=A_\pp$ is pure by hypothesis. Whence the conclusion.

\begin{proof}[Proof of the purity theorem]
Let us first show (i) by induction on the dimension. Let $A$ be a
noetherian local ring of \emph{dimension} $2$. Set $X'=\Spec(A)$,
$x=\rr(A)$, $X=X'-\{x\}$. One has $\prof A=2$. One can therefore apply
lemma \Ref{X.3.5} to the pair $(X', \{x\})$ and hence
$\Et(X')\to\Et(X)$ is fully faithful. Let now $r\colon R\to X$ be an
\'etale covering defined by a locally free coherent and \'etale
$\OX$-Algebra $\Aa=r_\ast(\Oo_R)$. Let us denote by $i\colon X\to X'$
the canonical immersion of $X$ in $X'$. I say that $i_\ast(\Aa)=\Bb$ is
a \emph{coherent} $\Oo_{X'}$-Algebra. Indeed, it suffices to apply the
``finiteness theorem'' \Ref{VIII}~\Ref{VIII.2.3}. I say that this
algebra is of depth $\geq 2$ at $x$. Indeed, it is the direct image of
an $\OX$-Module, with $X=X'-\{x\}$. Since $A$ is a \emph{regular} ring
of dimension $2$ one has: $\dimp \Bb + \prof\Bb=\dim A=2$, where
$\dimp\Bb$ denotes the projective dimension of $\Bb$. Hence
$\dimp\Bb=0$, hence $\Bb$ is projective, hence free. It follows that
$\Bb$ defines a finite flat covering of $X'=\Spec(A)$. The set of
points of $X'$ where this covering is not \'etale is a closed subset of
$X'$ whose equation is a principal ideal: the discriminant ideal of
$\Bb/A$. Now, by construction, this closed set is contained in
$x=\rr(A)$ hence is empty because $\dim A=2$.

Let $A$ be a regular noetherian local ring, $\dim A=n\geq3$. Suppose
(i) proved for rings of dimension $<n$. To prove that $A$ is pure, one
may assume $A$ complete by \Ref{X.3.8}. Let $t\in\rr(A)$ whose image in
$\rr(A)/\rr(A)^2$ is nonzero. Then $B=A/tA$ is a \emph{regular}
noetherian local ring of dimension $n-1$ hence is pure, because
$n-1\geq2$. One concludes by lemma \Ref{X.3.9}~(i), which is applicable
because $A$ is complete.

\emph{Let us show} (ii). Let $A$ be a noetherian local ring of
dimension $\geq 3$. Suppose there exists a \emph{regular} noetherian
local ring $B$ and a $B$-sequence $(t_1, \ldots, t_k)$ such that
$A\simeq B/(t_1, \ldots, t_k)$. Let us prove that $A$ is pure, by
induction on $k$. If $k=0$, one knows it by (i). \emph{Suppose
$k\geq1$ and the result acquired for $k'<k$.} By corollary
\Ref{X.3.9} one may assume that $A$ (hence also $B$) is complete. Set
$C=B/(t_1, \ldots, t_{k-1})$, hence $A\simeq C/t_kC$ and $t_k$ is
$C$-regular. By the induction hypothesis one knows that $C$ is pure, it
suffices to prove that lemma \Ref{X.3.9}~(ii) is applicable. Notation:
$A$ and $B$ of the lemma become $C$ and $A$. One has $\dim C\geq4$,
hence for every prime ideal $\pp$ of $C$ such that $\dim C/\pp=1$, one
has $\prof C_\pp\geq3$. Moreover, $C_\pp$ is a complete intersection
with $k'\leq k-1$, hence is pure by the induction hypothesis.
\end{proof}

\begin{theorem} \label{X.3.10}
Let $X$ be a locally noetherian prescheme and let $Y$ be a closed
subset of $X$. Suppose that one has $\Leff(X, Y)$ (\cf Examples
\Ref{X.2.1} and \Ref{X.2.2}). Suppose moreover that, for every open
neighbourhood $U$ of $Y$ and every $x\in X-U$, the local ring
$\Oo_{X, x}$ is regular of dimension $\geq2$ or else a complete
intersection of dimension $\geq3$. Then
\[
\pi_0(Y) \to \pi_0(X)
\]
is a bijection, and if $X$ is connected
\[
\pi_1(Y) \to \pi_1(X)
\]
is an isomorphism.
\end{theorem}

There is nothing left to prove. One remarks that, in the two examples
cited \Ref{X.2.1} and \Ref{X.2.2}, the complement of $U$ is a union of
finitely many closed points, whence it follows that the hypothesis on
the dimension of $\Oo_{X, x}$ is not a joke.

\numberwithin{equation}{section}
